% GENERATED FILE. Edit canonical lessons, then run npm run generate:books.
% canonical-source-hash: fnv1a64:f0352934ac516a3d
% canonical-lessons: RU-C163-sosiska, RU-C163-krevetka, RU-C163-borshch, RU-C163-shchi, RU-C163-blin

\chapter{Frankfurter, Prawn, Borscht, Cabbage soup, Pancake}
\label{ch:ru-frankfurter-prawn-borscht-cabbage-soup-pancake}
\hlchaptermodality{163}

\begin{chapteropening}
\textbf{By the end of this chapter:} \emph{I can say a frankfurter, a prawn, borscht, cabbage soup and a pancake.}

Complete the last lesson of chapter 163: I can say a frankfurter, a prawn, borscht, cabbage soup and a pancake.
\end{chapteropening}

\section[\texorpdfstring{\ru{сосиска} (sosíska)}{сосиска (sosíska)}]{\ru{сосиска} (sosíska) — a frankfurter}

\label{lesson:RU-C163-sosiska}



\begin{quote}
Before the new one: say the Russian for trainers, then the Russian for slippers.
\end{quote}

\subsection*{You'll want to know: \ru{сосиска}}
\textbf{\ru{сосиска}} (\emph{sosíska}) — \textquotedblleft{}a frankfurter\textquotedblright{}.

\begin{grammarlens}[title={What you've built}]
One more word for this chapter.
\end{grammarlens}

\subsection*{Your turn}
\begin{itemize}
\raggedright
  \item \emph{Say it:} \emph{sosíska}
  \item \emph{Say it:} \emph{sosíska}, once more
  \item \emph{Say it:} say \emph{sosíska}
  \item \emph{Recall:} say the Russian for trainers, then the Russian for slippers, then say \emph{sosíska} again
\end{itemize}

\begin{tcolorbox}[breakable,colback=teal!4,colframe=teal!35!black,title={Before you move on}]
What does \ru{сосиска} mean? (\textquotedblleft{}A frankfurter\textquotedblright{}.) Say it once more.
\end{tcolorbox}
\section[\texorpdfstring{\ru{креветка} (krevétka)}{креветка (krevétka)}]{\ru{креветка} (krevétka) — a prawn}

\label{lesson:RU-C163-krevetka}



\begin{quote}
Before the new one: say the Russian for slippers, then the Russian for a frankfurter.
\end{quote}

\subsection*{You'll want to know: \ru{креветка}}
\textbf{\ru{креветка}} (\emph{krevétka}) — \textquotedblleft{}a prawn\textquotedblright{}.

\begin{grammarlens}[title={What you've built}]
One more word for this chapter.
\end{grammarlens}

\subsection*{Your turn}
\begin{itemize}
\raggedright
  \item \emph{Say it:} \emph{krevétka}
  \item \emph{Say it:} \emph{krevétka}, once more
  \item \emph{Say it:} say \emph{krevétka}
  \item \emph{Recall:} say the Russian for slippers, then the Russian for a frankfurter, then say \emph{krevétka} again
\end{itemize}

\begin{tcolorbox}[breakable,colback=teal!4,colframe=teal!35!black,title={Before you move on}]
What does \ru{креветка} mean? (\textquotedblleft{}A prawn\textquotedblright{}.) Say it once more.
\end{tcolorbox}
\section[\texorpdfstring{\ru{борщ} (borshch)}{борщ (borshch)}]{\ru{борщ} (borshch) — borscht}

\label{lesson:RU-C163-borshch}



\begin{quote}
Before the new one: say the Russian for a frankfurter, then the Russian for a prawn.
\end{quote}

\subsection*{You'll want to know: \ru{борщ}}
\textbf{\ru{борщ}} (\emph{borshch}) — \textquotedblleft{}borscht\textquotedblright{}.

\begin{grammarlens}[title={What you've built}]
One more word for this chapter.
\end{grammarlens}

\subsection*{Your turn}
\begin{itemize}
\raggedright
  \item \emph{Say it:} \emph{borshch}
  \item \emph{Say it:} \emph{borshch}, once more
  \item \emph{Say it:} say \emph{borshch}
  \item \emph{Recall:} say the Russian for a frankfurter, then the Russian for a prawn, then say \emph{borshch} again
\end{itemize}

\begin{tcolorbox}[breakable,colback=teal!4,colframe=teal!35!black,title={Before you move on}]
What does \ru{борщ} mean? (\textquotedblleft{}Borscht\textquotedblright{}.) Say it once more.
\end{tcolorbox}
\section[\texorpdfstring{\ru{щи} (shchi)}{щи (shchi)}]{\ru{щи} (shchi) — cabbage soup}

\label{lesson:RU-C163-shchi}



\begin{quote}
Before the new one: say the Russian for a prawn, then the Russian for borscht.
\end{quote}

\subsection*{You'll want to know: \ru{щи}}
\textbf{\ru{щи}} (\emph{shchi}) — \textquotedblleft{}cabbage soup\textquotedblright{}.

\begin{grammarlens}[title={What you've built}]
One more word for this chapter.
\end{grammarlens}

\subsection*{Your turn}
\begin{itemize}
\raggedright
  \item \emph{Say it:} \emph{shchi}
  \item \emph{Say it:} \emph{shchi}, once more
  \item \emph{Say it:} say \emph{shchi}
  \item \emph{Recall:} say the Russian for a prawn, then the Russian for borscht, then say \emph{shchi} again
\end{itemize}

\begin{tcolorbox}[breakable,colback=teal!4,colframe=teal!35!black,title={Before you move on}]
What does \ru{щи} mean? (\textquotedblleft{}Cabbage soup\textquotedblright{}.) Say it once more.
\end{tcolorbox}
\section[\texorpdfstring{\ru{блин} (blin)}{блин (blin)}]{\ru{блин} (blin) — a pancake}

\label{lesson:RU-C163-blin}



\begin{quote}
Before the new one: say the Russian for borscht, then the Russian for cabbage soup.
\end{quote}

\subsection*{You'll want to know: \ru{блин}}
\textbf{\ru{блин}} (\emph{blin}) — \textquotedblleft{}a pancake\textquotedblright{}.

\begin{grammarlens}[title={What you've built}]
One more word for this chapter.
\end{grammarlens}

\subsection*{Your turn}
\begin{itemize}
\raggedright
  \item \emph{Say it:} \emph{blin}
  \item \emph{Say it:} \emph{blin}, once more
  \item \emph{Say it:} say \emph{blin}
  \item \emph{Recall:} say the Russian for borscht, then the Russian for cabbage soup, then say \emph{blin} again
\end{itemize}

\begin{tcolorbox}[breakable,colback=teal!4,colframe=teal!35!black,title={Before you move on}]
What does \ru{блин} mean? (\textquotedblleft{}A pancake\textquotedblright{}.) Say it once more.
\end{tcolorbox}
